\oldpage[0\textsubscript{I}]{34}

\subsection{Imágenes inversas de prehaces de grupos o de anillos}

\begin{env}[3.7.1]
\label{0.3.7.1-fr}
Vamos a mostrar que, cuando se toma como $\boldsymbol{K}$ la categoría de los
conjuntos, la imagen inversa por $\psi$ de todo prehaz $\mathcal{G}$ con
valores en $\boldsymbol{K}$ \emph{existe siempre} (con las notaciones y las
hipótesis sobre $X$, $Y$, $\psi$ de (3.5.3)). En efecto, para todo abierto
$U\subset X$, definamos $\mathcal{G}'(U)$ del modo siguiente. Un elemento $s'$
de $\mathcal{G}'(U)$ es una familia $(s'_x)_{x\in U}$, donde
$s'_x\in\mathcal{G}_{\psi(x)}$ para todo $x\in U$ y donde, para todo $x\in U$,
se cumple la condición siguiente: existen una vecindad abierta $V$ de
$\psi(x)$ en $Y$, una vecindad $W\subset\psi^{-1}(V)\cap U$ de $x$ y un
elemento $s\in\mathcal{G}(V)$ tales que $s'_z=s_{\psi(z)}$ para todo $z\in W$.
Se verifica inmediatamente que $U\to\mathcal{G}'(U)$ satisface los axiomas
de los \emph{haces}.

Sea ahora $\mathcal{F}$ un haz de conjuntos sobre $X$, y sean
$u:\mathcal{G}\to\psi_*(\mathcal{F})$,
$v:\mathcal{G}'\to\mathcal{F}$ morfismos. Se definen $u^\sharp$ y $v^\flat$
del modo siguiente. Si $s'$ es una sección de $\mathcal{G}'$ sobre una vecindad
$U$ de $x\in X$ y si $V$ es una vecindad abierta de $\psi(x)$ y
$s\in\mathcal{G}(V)$ es tal que $s'_z=s_{\psi(z)}$ para $z$ en una vecindad de
$x$ contenido en $\psi^{-1}(V)\cap U$, se toma
$u_x^\sharp(s'_x)=u_{\psi(x)}(s_{\psi(x)})$. Del mismo modo, si
$s\in\mathcal{G}(V)$ ($V$ abierto en $Y$), $v^\flat(s)$ es la sección de
$\mathcal{F}$ sobre $\psi^{-1}(V)$, imagen por $v$ de la sección $s'$ de
$\mathcal{G}'$ tal que $s'_x=s_{\psi(x)}$ para todo $x\in\psi^{-1}(V)$.
Además, el homomorfismo canónico (3.5.3)
$\rho:\mathcal{G}\to\psi_*(\psi^*(\mathcal{G}))$ se define del modo siguiente.
Para todo abierto $V\subset Y$ y toda sección $s\in\Gamma(V,\mathcal{G})$,
$\rho(s)$ es la sección $(s_{\psi(x)})_{x\in\psi^{-1}(V)}$ de
$\psi^*(\mathcal{G})$ sobre $\psi^{-1}(V)$. La comprobación de las relaciones
$(u^\sharp)^\flat=u$, $(v^\flat)^\sharp=v$ y $v^\flat=\psi_*(v)\circ\rho$
es inmediata y demuestra nuestra afirmación.

Se verifica que, si $w:\mathcal{G}_1\to\mathcal{G}_2$ es un homomorfismo de
prehaces de conjuntos sobre $Y$, $\psi^*(w)$ se describe explícitamente del
modo siguiente. Si $s'=(s'_x)_{x\in U}$ es una sección de
$\psi^*(\mathcal{G}_1)$ sobre un abierto $U$ de $X$, $(\psi^*(w))(s')$ es
la familia $(w_{\psi(x)}(s'_x))_{x\in U}$. Por último, es inmediato que para
todo abierto $V$ de $Y$, la imagen inversa de $\mathcal{G}|V$ por la
restricción de $\psi$ a $\psi^{-1}(V)$ es idéntica al haz inducido
$\psi^*(\mathcal{G})|\psi^{-1}(V)$.

Cuando $\psi$ es la identidad $1_X$, se recupera la definición de un haz de
conjuntos asociado a un prehaz (G, II, 1.2). Las consideraciones anteriores se
aplican sin cambios cuando $\boldsymbol{K}$ es la categoría de los grupos o
de los anillos (no necesariamente conmutativos).

Cuando $X$ es una parte cualquiera de un espacio topológico $Y$, y $j$ es la
inyección canónica $X\to Y$, para todo haz $\mathcal{G}$ sobre $Y$ con valores
en una categoría $\boldsymbol{K}$ se llama haz \emph{inducido} sobre $X$ por
$\mathcal{G}$ a la imagen inversa $j^*(\mathcal{G})$ (cuando existe). Para los
haces de conjuntos (o de grupos, o de anillos) se recupera la definición
usual (G, II, 1.5).
\end{env}
